\appendix
\chapter{Phụ lục}

\section{Danh sách tuyến API chính}

\begin{table}[H]
\centering
\scriptsize
\renewcommand{\arraystretch}{1.12}
\caption{Các tuyến API chính của hệ thống}
\begin{tabularx}{\textwidth}{L{4.8cm}L{2.2cm}Y}
\toprule
\textbf{Tuyến API} & \textbf{Thành phần} & \textbf{Vai trò} \\
\midrule
\code{GET} \code{/v1/ai-profile} & Medusa & Lấy hồ sơ AI hiện tại của người dùng. \\
\addlinespace
\code{POST} \code{/v1/ai-profile} & Medusa & Tạo hoặc cập nhật hồ sơ AI và ảnh cơ thể. \\
\addlinespace
\code{POST} \code{/v1/stylist/search} & Medusa & Nhận truy vấn từ giao diện, gọi dịch vụ AI và bổ sung dữ liệu cho phương án phối đồ. \\
\addlinespace
\code{GET} \code{/v1/stylist/sessions} & Medusa & Lấy lịch sử phiên tư vấn AI Stylist. \\
\addlinespace
\code{POST} \code{/v1/stylist/replace} & Medusa & Tìm sản phẩm tương tự để thay thế trong phương án phối đồ. \\
\addlinespace
\code{POST} \code{/api/v1/stylist/search} & FastAPI & Tuyến nội bộ xử lý luồng AI Stylist. \\
\addlinespace
\code{POST} \code{/api/v1/stylist/replace} & FastAPI & Tuyến nội bộ tìm sản phẩm tương đồng bằng vector. \\
\addlinespace
\code{GET} \code{/v1/user/assets} & Medusa & Lấy danh sách tủ đồ cá nhân. \\
\addlinespace
\code{POST} \code{/v1/user/garments} & Medusa & Tải trang phục cá nhân và phát sự kiện đồng bộ vector. \\
\addlinespace
\code{DELETE} \code{/v1/user/garments/[id]} & Medusa & Xóa trang phục cá nhân và phát sự kiện xóa vector. \\
\addlinespace
\code{POST} \code{/v1/user/garments/vton-result} & Medusa & Lưu ảnh kết quả VTON vào bộ sưu tập trang phục. \\
\addlinespace
\code{POST} \code{/v1/vton/jobs} & Medusa & Tạo tác vụ VTON và phát thông điệp vào \code{ai_vision_jobs}. \\
\addlinespace
\code{GET} \code{/v1/stream/vision-results} & Medusa & Gửi kết quả VTON tới giao diện bằng SSE. \\
\addlinespace
\code{POST} \code{/v1/checkout} & Medusa & Thực hiện quy trình thanh toán mô phỏng bằng Saga. \\
\addlinespace
\code{POST} \code{/v1/internal/products} & Medusa & Tạo sản phẩm từ công cụ quản trị nội bộ. \\
\bottomrule
\end{tabularx}
\end{table}

\section{Danh sách hàng đợi thông điệp}

\begin{table}[H]
\centering
\caption{Các hàng đợi RabbitMQ trong hệ thống}
\begin{tabularx}{\textwidth}{L{4cm}Y}
\toprule
\textbf{Hàng đợi} & \textbf{Vai trò} \\
\midrule
\code{product_metadata_sync} & Đồng bộ siêu dữ liệu sản phẩm từ Medusa sang LanceDB. \\
\addlinespace
\code{closet_metadata_sync} & Đồng bộ trang phục cá nhân khi thêm hoặc xóa khỏi tủ đồ. \\
\addlinespace
\code{ai_vision_jobs} & Chứa tác vụ VTON cần dịch vụ AI xử lý. \\
\addlinespace
\code{ai_vision_results} & Chứa kết quả VTON để Medusa cập nhật cơ sở dữ liệu và phát sự kiện SSE. \\
\bottomrule
\end{tabularx}
\end{table}

\section{Lệnh chạy hệ thống cục bộ}

\begin{enumerate}
    \item Cài đặt các gói phụ thuộc Node.js bằng lệnh \code{pnpm} \code{install}.
    \item Khởi động hạ tầng Docker bằng \code{docker} \code{compose} \code{up} \code{-d}.
    \item Khi PostgreSQL, Redis và RabbitMQ sẵn sàng, thực hiện chuyển đổi lược đồ Medusa bằng \code{pnpm} \code{-F} \code{@fundamental/api-medusa} \code{exec} \code{medusa} \code{db:migrate}.
    \item Khởi chạy giao diện và Medusa từ thư mục gốc bằng \code{pnpm} \code{dev}.
    \item Chạy dịch vụ AI trong \code{apps/ai-service} bằng \code{uv} \code{run} \code{uvicorn} \code{main:app} \code{--reload}.
\end{enumerate}

\clearpage
